\documentclass[10pt,a4paper]{amsart}
\usepackage{amsmath,amssymb,amsthm,graphicx}
\usepackage{verbatim}
\newenvironment{reference}{\comment}{\endcomment}
%\newenvironment{reference}{}{}
\newenvironment{slogan}{\comment}{\endcomment}
\newenvironment{history}{\comment}{\endcomment}

% For commutative diagrams we use Xy-pic
\usepackage[all]{xy}

% We use 2cell for 2-commutative diagrams.
\xyoption{2cell}
% Theorem environments.
%
\theoremstyle{plain}
\newtheorem{theorem}[subsection]{Theorem}
\newtheorem{proposition}[subsection]{Proposition}
\newtheorem{lemma}[subsection]{Lemma}

\theoremstyle{definition}
\newtheorem{definition}[subsection]{Definition}
\newtheorem{example}[subsection]{Example}
\newtheorem{exercise}[subsection]{Exercise}
\newtheorem{situation}[subsection]{Situation}

\theoremstyle{remark}
\newtheorem{remark}[subsection]{Remark}
\newtheorem{remarks}[subsection]{Remarks}

\numberwithin{equation}{subsection}

% Macros
%
\def\lim{\mathop{\mathrm{lim}}\nolimits}
\def\colim{\mathop{\mathrm{colim}}\nolimits}
\def\Spec{\mathop{\mathrm{Spec}}}
\def\Hom{\mathop{\mathrm{Hom}}\nolimits}
\def\Ext{\mathop{\mathrm{Ext}}\nolimits}
\def\SheafHom{\mathop{\mathcal{H}\!\mathit{om}}\nolimits}
\def\SheafExt{\mathop{\mathcal{E}\!\mathit{xt}}\nolimits}
\def\Sch{\mathit{Sch}}
\def\Mor{\mathop{\mathrm{Mor}}\nolimits}
\def\Ob{\mathop{\mathrm{Ob}}\nolimits}
\def\Sh{\mathop{\mathit{Sh}}\nolimits}
\def\NL{\mathop{N\!L}\nolimits}
\def\CH{\mathop{\mathrm{CH}}\nolimits}
\def\proetale{{pro\text{-}\acute{e}tale}}
\def\etale{{\acute{e}tale}}
\def\QCoh{\mathit{QCoh}}
\def\Ker{\mathop{\mathrm{Ker}}}
\def\Im{\mathop{\mathrm{Im}}}
\def\Coker{\mathop{\mathrm{Coker}}}
\def\Coim{\mathop{\mathrm{Coim}}}

% Boxtimes
%
\DeclareMathSymbol{\boxtimes}{\mathbin}{AMSa}{"02}

%
% Macros for moduli stacks/spaces
%
\def\QCohstack{\mathcal{QC}\!\mathit{oh}}
\def\Cohstack{\mathcal{C}\!\mathit{oh}}
\def\Spacesstack{\mathcal{S}\!\mathit{paces}}
\def\Quotfunctor{\mathrm{Quot}}
\def\Hilbfunctor{\mathrm{Hilb}}
\def\Curvesstack{\mathcal{C}\!\mathit{urves}}
\def\Polarizedstack{\mathcal{P}\!\mathit{olarized}}
\def\Complexesstack{\mathcal{C}\!\mathit{omplexes}}
% \Pic is the operator that assigns to X its picard group, usage \Pic(X)
% \Picardstack_{X/B} denotes the Picard stack of X over B
% \Picardfunctor_{X/B} denotes the Picard functor of X over B
\def\Pic{\mathop{\mathrm{Pic}}\nolimits}
\def\Picardstack{\mathcal{P}\!\mathit{ic}}
\def\Picardfunctor{\mathrm{Pic}}
\def\Deformationcategory{\mathcal{D}\!\mathit{ef}}
\usepackage{cleveref}
\crefname{originallemma}{original lemma}{original lemmas}
\Crefname{originallemma}{Original lemma}{Original lemmas}
\makeatletter
\newlabel{groupoids-lemma-group-scheme-over-field-separated}{{047L}{source}{}{sourcepointer-1}{}}
\newlabel{groupoids-lemma-group-scheme-over-field-separated@cref}{{[originallemma][0][]047L}{[1][0][]source}{}{}{}}
\makeatother
\crefname{originalsection}{original section}{original sections}
\Crefname{originalsection}{Original section}{Original sections}
\makeatletter
\newlabel{schemes-example-affine-space-zero-doubled}{{01JD}{source}{}{sourcepointer-2}{}}
\newlabel{schemes-example-affine-space-zero-doubled@cref}{{[originalsection][0][]01JD}{[1][0][]source}{}{}{}}
\makeatother
\title{A flat group scheme of finite type over the affine line which is not separated}
\author{The Stacks Project Authors}
\date{stacks-project:examples.tex:local-snapshot-20261012; GNU FDL 1.2 or later}
\begin{document}
\maketitle
\noindent\textit{Editorial note.} All delivered mathematical text is the author's own, copied from the pinned Stacks Project chapter file examples.tex; no statement or proof step was rewritten, reordered or removed.\par\medskip
\noindent\textit{Editorial note.} The counted proof is the author's complete argument: The counted proof is 'See the discussion above' - the construction is the context span, which is the same material delivered as the context of L48293; content verbatim.\par\medskip
\noindent\textit{Editorial note.} The label lemma-non-separated-group-scheme names the counted statement, and the proof environment immediately following it in the source is delivered verbatim as its proof.\par\medskip
\noindent\textit{Editorial note.} The counted statement is carried by the author's own lemma environment in the chapter source.\par\medskip
\noindent\textit{Allowed standard prerequisites (editorial).}\begin{itemize}
\item Context span 3830-3856: G = affine line with 0 doubled over S = A\textasciicircum{}1\_k, fibres singletons or two-point sets with the U-point as zero, opens U, V mapping isomorphically to S with U cap V -> S \textbackslash{} \{0\}, multiplication defined piecewise on the four pieces, author's own omitted group-structure verification
\item Schemes, example-affine-space-zero-doubled; the diagonal criterion for separatedness
\end{itemize}\medskip
\ifdefined\hypertarget\hypertarget{sourcepointer-1}{}\fi
\noindent\textit{Source reference.} Original lemma 047L: sources/stacks-chow/tags.txt, source lines 5512--5512. Prerequisite pointer invoked at 3871; it names a standard fact of a sibling chapter that the counted argument uses, and every step of the counted proof built on it is retained verbatim in the delivered body.\par\medskip
\ifdefined\hypertarget\hypertarget{sourcepointer-2}{}\fi
\noindent\textit{Source reference.} Original section 01JD: sources/stacks-chow/tags.txt, source lines 2010--2010. Prerequisite pointer invoked at 3877; it names a standard fact of a sibling chapter that the counted argument uses, and every step of the counted proof built on it is retained verbatim in the delivered body.\par\medskip
\label{section-non-separated-group-scheme}

\noindent
Every group scheme over a field is separated, see
Groupoids, Lemma \ref{groupoids-lemma-group-scheme-over-field-separated}.
This is not true for group schemes over a base.

\medskip\noindent
Let $k$ be a field. Let $S = \Spec(k[x]) = \mathbf{A}^1_k$.
Let $G$ be the affine line with $0$ doubled (see
Schemes, Example \ref{schemes-example-affine-space-zero-doubled})
seen as a scheme over $S$. Thus a fibre of $G \to S$ is either a
singleton or a set with two elements (one in $U$ and one in $V$).
Thus we can endow these fibres with the structure of a group (by
letting the element in $U$ be the zero of the group structure).
More precisely, $G$ has two opens $U, V$ which map isomorphically
to $S$ such that $U \cap V$ is mapped isomorphically to
$S \setminus \{0\}$. Then
$$
G \times_S G = U \times_S U \cup V \times_S U \cup
U \times_S V \cup V \times_S V
$$
where each piece is isomorphic to $S$. Hence we can define a multiplication
$m : G \times_S G \to G$ as the unique $S$-morphism which maps the first
and the last piece into $U$ and the two middle pieces into $V$. This matches
the pointwise description given above. We omit the
verification that this defines a group scheme structure.

\begin{lemma}
\label{lemma-non-separated-group-scheme}
There exists a flat group scheme of finite type over the affine line
which is not separated.
\end{lemma}

\begin{proof}
See the discussion above.
\end{proof}

\end{document}
